\documentclass[12pt]{article}%
\usepackage{amsfonts}
\usepackage{fancyhdr}
\usepackage{comment}
\usepackage[a4paper, top=2.5cm, bottom=2.5cm, left=2.2cm, right=2.2cm]%
{geometry}
\usepackage{times}
\usepackage{amsmath}
\usepackage{changepage}
\usepackage{amssymb}
\usepackage{graphicx}

\setcounter{MaxMatrixCols}{30}
\newtheorem{theorem}{Theorem}
\newtheorem{acknowledgement}[theorem]{Acknowledgement}
\newtheorem{algorithm}[theorem]{Algorithm}
\newtheorem{axiom}{Axiom}
\newtheorem{case}[theorem]{Case}
\newtheorem{claim}[theorem]{Claim}
\newtheorem{conclusion}[theorem]{Conclusion}
\newtheorem{condition}[theorem]{Condition}
\newtheorem{conjecture}[theorem]{Conjecture}
\newtheorem{corollary}[theorem]{Corollary}
\newtheorem{criterion}[theorem]{Criterion}
\newtheorem{definition}[theorem]{Definition}
\newtheorem{example}[theorem]{Example}
\newtheorem{exercise}[theorem]{Exercise}
\newtheorem{lemma}[theorem]{Lemma}
\newtheorem{notation}[theorem]{Notation}
\newtheorem{problem}[theorem]{Problem}
\newtheorem{proposition}[theorem]{Proposition}
\newtheorem{remark}[theorem]{Remark}
\newtheorem{solution}[theorem]{Solution}
\newtheorem{summary}[theorem]{Summary}
\newenvironment{proof}[1][Proof]{\textbf{#1.} }{\ \rule{0.5em}{0.5em}}

\newcommand{\Q}{\mathbb{Q}}
\newcommand{\R}{\mathbb{R}}
\newcommand{\C}{\mathbb{C}}
\newcommand{\Z}{\mathbb{Z}}

%%%%
% ME %
%%%%

\usepackage{listings}
\usepackage{color}
\usepackage{xcolor}
\usepackage{mdframed}

\definecolor{dkgreen}{rgb}{0,0.6,0}
\definecolor{gray}{rgb}{0.5,0.5,0.5}
\definecolor{mauve}{rgb}{0.58,0,0.82}

\lstset{%frame=tb,
language=Matlab,
aboveskip=3mm,
belowskip=3mm,
showstringspaces=false,
columns=flexible,
basicstyle={\small\ttfamily},
numbers=none,
numberstyle=\tiny\color{gray},
keywordstyle=\color{blue},
commentstyle=\color{dkgreen},
stringstyle=\color{mauve},
breaklines=true,
breakatwhitespace=true,
tabsize=3
}

\usepackage{outlines}
\usepackage{enumitem}
%\setenumerate[1]{label=\Roman*.}
%\setenumerate[2]{label=\Alph*.}
%\setenumerate[3]{label=\roman*.}
%\setenumerate[4]{label=\alph*.}

\usepackage{titlesec} 

\titleformat{\subsection}[runin]
  {\normalfont\large\bfseries}{\thesubsection}{1em}{}
\titleformat{\subsubsection}[runin]
  {\normalfont\normalsize\bfseries}{\thesubsubsection}{1em}{}

% for enumerate spacing
\newenvironment{tight_enum}{
\begin{enumerate}[label=\alph*.]
\setlength{\itemsep}{0pt}
\setlength{\parskip}{0pt}
}{\end{enumerate}}

\usepackage{exercise}

\setlength\parindent{0pt}

\newcommand{\code}[1]{\texttt{#1}}
\newcommand{\shp}{$>>>$ }
\newcommand{\tab}{\hspace*{45pt}}

\usepackage{multicol}
\setlength{\columnsep}{0.6cm}

\begin{document}

\title{Homework 08: Chapter 10}

\author{EEC 3150}
\date{}

\maketitle

\begin{center}
\fbox{\fbox{\parbox{5.5in}{\centering
For each Python exercise, write a separate Python script. Submit all scripts on Blackboard under the appropriate assignment. Name your script files with the following format: \\
\vspace{10pt}
HW$<$assignment number$>$\_Ex$<$exercise number$>$\_$<$FirstInitial$><$LastName$>$.py \\
\vspace{10pt}
Example: HW01\_Ex01\_GWest.py \\
\vspace{10pt}
For short answer or math exercises, either do your work on a sheet of paper and upload a picture of it to Blackboard or submit some kind of text file or Word document to Blackboard. You can place all such exercises in a single document as long as you clearly label each exercise. These files should follow a similar naming convention as the Python scripts.
}}}
\end{center}


\pagebreak


\begin{multicols*}{2}

%%% NEW EXERCISE %%%
\begin{Exercise}[origin={10 points, Python}]

Create a class Complex which implements complex number arithmetic:

\begin{enumerate}[leftmargin=*]
	\item The class should have two attributes:
	\code{
	\begin{itemize}[leftmargin=*]
		\item \_real
		\item \_imag
	\end{itemize}
	}
	\item The class should have these magic methods:
	\begin{itemize}[leftmargin=*]
		\item \code{\_\_str\_\_} should return a string in the following format: \code{a+bi}
		\item \code{\_\_add\_\_} should return a new Complex object that is the sum of the two operands
		\item \code{\_\_sub\_\_} should return a new Complex object that is the difference of the two operands
		\item \code{\_\_mul\_\_} should return a new Complex object that is the product of the two operands
		\item \code{\_\_truediv\_\_} should return a new Complex object that is the quotient of the two operands
	\end{itemize}
	\item The class should have one additional method:
	\begin{itemize}
		\item \code{conj} (complex conjugate)
	\end{itemize}
\end{enumerate}

\rule[0pt]{\linewidth}{1pt}

\vspace{4pt}

Here's some examples the class and methods:

\rule[0pt]{\linewidth}{1pt}

\vspace{10pt}

For reference, here are some definitions for copmlex arithmetic. Let $x = a + bi$ and $y = c + di$ be two complex numbers. \\

The sum of $x$ and $y$ is:
\begin{align*}
x + y &= (a + bi) + (c + di) \\
      &= (a + c) + (b + d)i
\end{align*}

The difference of $x$ and $y$ is:
\begin{align*}
x - y &= (a + bi) - (c + di) \\
      &= (a - c) + (b - d)i
\end{align*}

The product of $x$ and $y$ is:
\begin{align*}
xy &= (a + bi) (c + di) \\
   &= (ac - bd) + (ad + bc)i
\end{align*}

The quotient of $x$ and $y$ is:
\begin{align*}
\frac{x}{y} &= \frac{a + bi}{c + di} \\
            &= \frac{(ac + bd)}{c^2 + d^2} + \frac{(bc - ad)}{c^2 + d^2}i
\end{align*}

\end{Exercise}



%\vfill\null
%\columnbreak

%%% NEW EXERCISE %%%
\begin{Exercise}[origin={20 points, Python}]

\begin{enumerate}[leftmargin=*]
	\item Write two functions that perform the same task in different ways. The task should take anywhere between (roughly) 0.1 and 1 seconds to complete (basically it needs to take a while).
	\item Use the \code{time} module to find the average time it takes to perform each function. Clearly print average times, clearly label each version.
\end{enumerate}

\end{Exercise}



\vfill\null
\end{multicols*}




\end{document}






